\begin{abstract}
Adaptive testing with item response theory needs item parameters calibrated on other
systems, and new, private or costly benchmarks for language models often have only a
handful. We study this calibration uncertainty with EOST-M, which uses each item's expected
response function under an empirical-Bayes prior and selects items by expected information
gain. This model is exact given the calibrators' abilities, and a difficulty with Gaussian
uncertainty $s$ caps an item's Fisher information at $2/(\pi s^2)$ for any link. Two
frozen-plan studies on eight benchmarks give three messages.
First, plug-in Fisher-information testing fails with small panels: with eight calibrators and
one target per family cluster, EOST-M predicted held-out verdicts better (median
\STwoRTwoMed\%, significant on \STwoRTwoFav\ of 8 benchmarks). Second, accounting for the
uncertainty fixes much of this: EOST-M beat a regularized Rasch-type test (\STwoCOneMed\%),
\AbstractFB\ and its score error was lower than with random items under the same model,
ATLAS-style testing, tinyBenchmarks and MINCE, all run on the same eight calibrators, on every
benchmark, meeting the frozen criterion only against MINCE. Third, the information criterion then matters little.
EOST-M's 90\% intervals still under-cover (\STwoCoveigmg; \STwoCoveigmgRR\ after
recalibration).
\end{abstract}
